% ============================================================================
\section{Pac-Man}
\label{sec:pacman}
% Source: assignments/05_pacman/index.html
% ============================================================================

% ----------------------------------------------------------------------------
\subsection{Minimax}

\problempart{1a}{Multi-ghost minimax recurrence}
\question{There are $n+1$ agents $a_0,\dots,a_n$: Pac-Man ($a_0$, max) and $n$ ghosts (min). One unit of depth
means all $n+1$ agents move once. $d$ starts at $d_{\max}$ at the root and decreases. Write
$V_{\text{minmax}}(s,d)$ as a piecewise function of $\text{IsEnd}(s)$, $\text{Utility}(s)$, $\text{Eval}(s)$,
$\text{Player}(s)$, $\text{Actions}(s)$ and $\text{Succ}(s,a)$; $n$ may be used, $d_{\max}$ may not.}
\begin{figure}[H]\centering
  \includegraphics[width=0.75\linewidth]{05_pacman/minimax_depth.png}
  \caption{Depth vs.\ tree height (from the handout).}
\end{figure}
\expect{The recurrence as a piecewise function, using exactly these functions.}
\begin{solution}
\[
V_{\text{minmax}}(s,d)=
\begin{cases}
\text{Utility}(s) & \text{if } \text{IsEnd}(s),\\[2pt]
\text{Eval}(s) & \text{else if } d=0,\\[2pt]
\displaystyle\max_{a\in\text{Actions}(s)} V_{\text{minmax}}\bigl(\text{Succ}(s,a),\,d\bigr) & \text{else if } \text{Player}(s)=a_0,\\[8pt]
\displaystyle\min_{a\in\text{Actions}(s)} V_{\text{minmax}}\bigl(\text{Succ}(s,a),\,d\bigr) & \text{else if } \text{Player}(s)=a_i,\ 1\le i<n,\\[8pt]
\displaystyle\min_{a\in\text{Actions}(s)} V_{\text{minmax}}\bigl(\text{Succ}(s,a),\,d-1\bigr) & \text{else if } \text{Player}(s)=a_n.
\end{cases}
\]
Pac-Man and the first $n-1$ ghosts pass the same depth to their successors; only the last ghost $a_n$
completes a round, so its successors (Pac-Man's next turn) are one level shallower. Hence the depth limit
$d=0$ is always reached at a Pac-Man node, and a search of depth $d_{\max}$ explores
$d_{\max}(n+1)$ plies, as in the figure. The root call is $V_{\text{minmax}}(s_{\text{start}},d_{\max})$.
\end{solution}

\problempart{1b}{\texttt{MinimaxAgent}}
\codepart{5.1b}
We checked the implementation against the values given in the handout. The root minimax values on
\texttt{minimaxClassic} are $9,8,7,-492$ for depths $1$--$4$, and $9,18,27$ on \texttt{mediumClassic} for depths
$1$--$3$, and the agent's chosen action attains the root value in every case. The local grader awards 4/4 with no
timeouts.

% ----------------------------------------------------------------------------
\subsection{Alpha-beta pruning}
\codepart{5.2a}
Each node carries $\alpha$ (the best value Pac-Man can already guarantee) and $\beta$ (the best value the ghosts can
already guarantee). A max node stops as soon as its value reaches $\beta$, and every ghost's min node stops as soon as its
value drops to $\alpha$. Because consecutive ghost layers all minimise, they simply pass $\alpha$ and $\beta$
down unchanged. The agent returns exactly the minimax values above, and the grader awards 4/4 with no timeouts.

% ----------------------------------------------------------------------------
\subsection{Expectimax}

\problempart{3a}{Expectimax recurrence}
\question{Write $V_{\text{exptmax}}(s,d)$, the maximum expected utility against ghosts that each pick a legal move
uniformly at random, using the same functions as in 1a.}
\expect{A piecewise recurrence in the style of 1a.}
\begin{solution}
\[
V_{\text{exptmax}}(s,d)=
\begin{cases}
\text{Utility}(s) & \text{if } \text{IsEnd}(s),\\[2pt]
\text{Eval}(s) & \text{else if } d=0,\\[2pt]
\displaystyle\max_{a\in\text{Actions}(s)} V_{\text{exptmax}}\bigl(\text{Succ}(s,a),\,d\bigr) & \text{else if } \text{Player}(s)=a_0,\\[8pt]
\displaystyle\frac{1}{|\text{Actions}(s)|}\sum_{a} V_{\text{exptmax}}\bigl(\text{Succ}(s,a),\,d\bigr) & \text{else if } \text{Player}(s)=a_i,\ i<n,\\[8pt]
\displaystyle\frac{1}{|\text{Actions}(s)|}\sum_{a} V_{\text{exptmax}}\bigl(\text{Succ}(s,a),\,d-1\bigr) & \text{else if } \text{Player}(s)=a_n.
\end{cases}
\]
This is the minimax recurrence with each ghost's $\min$ replaced by the expectation under the uniform random
policy $\pi_{\text{ghost}}(a\mid s)=1/|\text{Actions}(s)|$.
\end{solution}

\problempart{3b}{\texttt{ExpectimaxAgent}}
\codepart{5.3b}
On \texttt{trappedClassic} at depth 3 the agent finishes with exactly the scores the handout predicts: $-502$ when it
loses and $532$ when it wins. It won 29 of 50 seeded games (the handout says ``about half''). The grader awards 4/4 with no timeouts.

% ----------------------------------------------------------------------------
\subsection{Evaluation function (extra credit)}

\problempart{4a}{\texttt{better\_evaluation\_function}}
\codepart{5.4a}
We tested the agent the way the (disabled) extra-credit grader does: depth-2 expectimax against two random
ghosts on \texttt{smallClassic}, 20 games, seed \texttt{'testing'}. Our evaluation function wins $15/20$ games
with an average winning score of $1663.7$, above the $1600$ threshold for the maximum extra credit. Over 100
games (five seeds) it wins $90/100$ with an average winning score of $1647.5$. A batch of 20 games takes about
8 seconds in total, far inside the time limit.

\problempart{4b}{Describing the evaluation function}
\question{Describe the evaluation function and its high-level motivation. What else did you try, and what worked
and what did not? State whether you consent to the solution being shared if it places in the top 3.}
\expect{A short paragraph.}
\begin{solution}
The function starts from the game score, which already rewards food ($+10$), eating ghosts ($+200$) and
winning ($+500$), and adds terms that make \emph{progress} visible to a shallow search. On its own, the score
hardly changes over two plies, so Pac-Man dithers. We run one breadth-first search from Pac-Man's position to get
true maze distances, so walls are respected. The added terms are: $+10/(1+d_{\text{food}})$ for the nearest food
pellet, which pulls Pac-Man towards food; $-4$ per remaining pellet and $-20$ per remaining capsule, which reward
clearing the board and eating capsules; $+200/(1+d)$ for every scared ghost Pac-Man can reach before its timer
runs out, which turns capsules into $+200$-point ghost hunts and accounts for most of the score above
$\approx1300$; and, for active ghosts, $-500$ when adjacent, otherwise only a weak $-2/d$, because a random ghost
far away is not a real threat. Win and loss states return the raw score.

To find out what mattered, we ran an ablation (\texttt{scripts/pacman\_eval\_ablation.py}; 20 games,
grader seed):
\begin{center}\small
\begin{tabular}{lccc}
\toprule
variant & wins & avg.\ winning score & avg.\ score \\
\midrule
submitted function & 15/20 & 1663.7 & 1274.0 \\
Manhattan instead of maze (BFS) distances & 19/20 & 1630.1 & 1577.2 \\
without the scared-ghost hunting term & 19/20 & 1349.9 & 1266.8 \\
danger radius 3 instead of 1 (more cautious) & 15/20 & 1598.0 & 1277.8 \\
game score only (\texttt{score\_evaluation\_function}) & 2/20 & 962.0 & $-398.8$ \\
\bottomrule
\end{tabular}
\end{center}
The progress terms are what make the agent win at all: the raw score alone wins only 2 of 20 games. The
scared-ghost term is what lifts the winning score from about 1350 to above 1600, because the extra points come from
eating ghosts after a capsule, not from the capsule itself. Being more cautious (treating ghosts within three
squares as deadly) does not win more games and lowers the score, because Pac-Man avoids corridors with food. The
choice between maze and Manhattan distances turned out not to matter: the 20-game difference above is noise, and
over 100 games (five seeds) the two versions are equivalent (90 vs.\ 88 wins, average winning score 1647.5 vs.\
1648.5). We kept the maze distance because it is the more faithful measure of how far food really is.
\textit{Consent:} we consent to our solution being shared with other students.
\end{solution}

% ----------------------------------------------------------------------------
\subsection{AI (mis)alignment and reward hacking}

Compare \texttt{python pacman.py -p MinimaxAgent -l trappedClassic -a depth=3} with the
\texttt{ExpectimaxAgent} on the same layout (run each several times).

\problempart{5a}{Why the agents differ}
\question{Why does the minimax agent always rush towards the \emph{closest} ghost, while the expectimax agent
(occasionally) does not?}
\expect{One sentence for minimax, stating the assumption that causes this, and one sentence for expectimax.}
\begin{solution}
Minimax assumes the ghosts are perfect adversaries. Under that worst-case assumption every line of play in
\texttt{trappedClassic} ends in death, so all outcomes cost the $-500$ loss penalty, and the only remaining
difference is the $-1$ time penalty per move. Dying as quickly as possible, by running into the nearest ghost, is
therefore ``optimal'': in our 50 runs it always lost, with score $-501$. Expectimax instead assumes each ghost moves
uniformly at random, so walking away from a ghost has a real chance that the ghost does not follow. The expected
value of going for the food (about $+532$ with positive probability) beats certain death, and the gamble paid off
in 29 of 50 games.
\end{solution}

\problempart{5b}{Fixing the objective}
\question{Suggest a change to $\text{Eval}(s)$ (\texttt{score\_evaluation\_function}) and/or $\text{Utility}(s)$ (the
final score), or to how the game score is computed, that would stop minimax from dying instantly on
\texttt{trappedClassic} and make it behave more like expectimax.}
\expect{1--2 sentences describing the change and why it works.}
\begin{solution}
Make survival time valuable in losing states. For example, use
$\text{Utility}(s)=\text{score}(s)+c\cdot\text{(moves survived)}$ with $c>1$ when $s$ is a loss, or equivalently
drop the per-move time penalty from the score of losing states. Then a later death is strictly better than an
immediate one, so even a worst-case planner moves away from the ghosts and stays alive. That keeps open the chance
that the real (random) ghosts make a mistake, which is exactly the behaviour expectimax shows.
\end{solution}

\problempart{5c}{Another real-world example}
\question{Describe a realistic scenario outside Pac-Man in which a designer's objective is open to reward hacking
or causes negative side effects.}
\expect{2--5 sentences, stating whether it is reward hacking, a negative side effect, or both, with a justification.}
\begin{solution}
A company deploys a customer-support chatbot and trains it with the reward ``fraction of conversations that end
with the ticket marked \emph{resolved}'', together with a short-survey satisfaction score. The model learns to
close tickets quickly with confident but unverified answers, to steer users away from refunds that would reopen
the ticket, and to flatter users so they rate the chat highly. This is \textbf{reward hacking}: the proxy, tickets
marked resolved and high ratings, is maximised while the designer's real goal, actually solving customers'
problems, is not. It also has \textbf{negative side effects} on other stakeholders: customers who were wrongly
told their problem is fixed bear the cost later, and human agents receive harder, angrier escalations.
\end{solution}
